\documentclass[10pt,a4paper]{amsart}
\usepackage{amsmath,amssymb,amsthm,graphicx}
% For commutative diagrams we use Xy-pic
\usepackage[all]{xy}

% We use 2cell for 2-commutative diagrams.
\xyoption{2cell}
% Theorem environments.
%
\theoremstyle{plain}
\newtheorem{theorem}[subsection]{Theorem}
\newtheorem{proposition}[subsection]{Proposition}
\newtheorem{lemma}[subsection]{Lemma}

\theoremstyle{definition}
\newtheorem{definition}[subsection]{Definition}
\newtheorem{example}[subsection]{Example}
\newtheorem{exercise}[subsection]{Exercise}
\newtheorem{situation}[subsection]{Situation}

\theoremstyle{remark}
\newtheorem{remark}[subsection]{Remark}
\newtheorem{remarks}[subsection]{Remarks}

\numberwithin{equation}{subsection}

% Macros
%
\def\lim{\mathop{\mathrm{lim}}\nolimits}
\def\colim{\mathop{\mathrm{colim}}\nolimits}
\def\Spec{\mathop{\mathrm{Spec}}}
\def\Hom{\mathop{\mathrm{Hom}}\nolimits}
\def\Ext{\mathop{\mathrm{Ext}}\nolimits}
\def\SheafHom{\mathop{\mathcal{H}\!\mathit{om}}\nolimits}
\def\SheafExt{\mathop{\mathcal{E}\!\mathit{xt}}\nolimits}
\def\Sch{\mathit{Sch}}
\def\Mor{\mathop{\mathrm{Mor}}\nolimits}
\def\Ob{\mathop{\mathrm{Ob}}\nolimits}
\def\Sh{\mathop{\mathit{Sh}}\nolimits}
\def\NL{\mathop{N\!L}\nolimits}
\def\CH{\mathop{\mathrm{CH}}\nolimits}
\def\proetale{{pro\text{-}\acute{e}tale}}
\def\etale{{\acute{e}tale}}
\def\QCoh{\mathit{QCoh}}
\def\Ker{\mathop{\mathrm{Ker}}}
\def\Im{\mathop{\mathrm{Im}}}
\def\Coker{\mathop{\mathrm{Coker}}}
\def\Coim{\mathop{\mathrm{Coim}}}

% Boxtimes
%
\DeclareMathSymbol{\boxtimes}{\mathbin}{AMSa}{"02}

%
% Macros for moduli stacks/spaces
%
\def\QCohstack{\mathcal{QC}\!\mathit{oh}}
\def\Cohstack{\mathcal{C}\!\mathit{oh}}
\def\Spacesstack{\mathcal{S}\!\mathit{paces}}
\def\Quotfunctor{\mathrm{Quot}}
\def\Hilbfunctor{\mathrm{Hilb}}
\def\Curvesstack{\mathcal{C}\!\mathit{urves}}
\def\Polarizedstack{\mathcal{P}\!\mathit{olarized}}
\def\Complexesstack{\mathcal{C}\!\mathit{omplexes}}
% \Pic is the operator that assigns to X its picard group, usage \Pic(X)
% \Picardstack_{X/B} denotes the Picard stack of X over B
% \Picardfunctor_{X/B} denotes the Picard functor of X over B
\def\Pic{\mathop{\mathrm{Pic}}\nolimits}
\def\Picardstack{\mathcal{P}\!\mathit{ic}}
\def\Picardfunctor{\mathrm{Pic}}
\def\Deformationcategory{\mathcal{D}\!\mathit{ef}}
\usepackage{cleveref}
\crefname{originallemma}{original lemma}{original lemmas}
\Crefname{originallemma}{Original lemma}{Original lemmas}
\makeatletter
\newlabel{topology-lemma-quasi-compact-open-immersion-constructible-image}{{09YD}{source}{}{sourcepointer-1}{}}
\newlabel{topology-lemma-quasi-compact-open-immersion-constructible-image@cref}{{[originallemma][0][]09YD}{[1][0][]source}{}{}{}}
\makeatother
\crefname{originallemma}{original lemma}{original lemmas}
\Crefname{originallemma}{Original lemma}{Original lemmas}
\makeatletter
\newlabel{topology-lemma-constructible-partition-refined-by-stratification}{{09Y4}{source}{}{sourcepointer-2}{}}
\newlabel{topology-lemma-constructible-partition-refined-by-stratification@cref}{{[originallemma][0][]09Y4}{[1][0][]source}{}{}{}}
\makeatother
\crefname{originallemma}{original lemma}{original lemmas}
\Crefname{originallemma}{Original lemma}{Original lemmas}
\makeatletter
\newlabel{properties-lemma-locally-constructible}{{054C}{source}{}{sourcepointer-3}{}}
\newlabel{properties-lemma-locally-constructible@cref}{{[originallemma][0][]054C}{[1][0][]source}{}{}{}}
\makeatother
\crefname{originaldefinition}{original definition}{original definitions}
\Crefname{originaldefinition}{Original definition}{Original definitions}
\makeatletter
\newlabel{definition-constructible}{{03RW}{source}{}{sourcepointer-4}{}}
\newlabel{definition-constructible@cref}{{[originaldefinition][0][]03RW}{[1][0][]source}{}{}{}}
\makeatother
\title{Characterisations of constructible sheaves on a quasi-compact quasi-separated scheme}
\author{The Stacks Project Authors}
\date{stacks-project:etale-cohomology.tex:local-snapshot-20261011; GNU FDL 1.2 or later}
\begin{document}
\maketitle
\noindent\textit{Editorial note.} All delivered mathematical text is the author's own, copied from the pinned Stacks Project chapter file etale-cohomology.tex; no statement or proof step was rewritten, reordered or removed.\par\medskip
\noindent\textit{Editorial note.} The proof is the author's complete argument for the equivalence, including the refinement of the partitions and the remark that the same proof works for abelian sheaves and modules over a Noetherian ring.\par\medskip
\noindent\textit{Editorial note.} Every labelled reference occurring inside the counted spans is carried in assembly.external\_references as a non-load-bearing navigation pointer (same-chapter ones at their definition site, cross-chapter ones at the invocation inside the counted span), so no load-bearing mathematics is replaced by a stub. No bibliography citation occurs inside the counted spans.\par\medskip
\noindent\textit{Editorial note.} The checked chapter list for same-name collisions is the set of stacks chapters staged in this repository: chow, cohomology, divisors, etale-cohomology, flat, more-algebra, more-morphisms, morphisms, spaces-more-morphisms, varieties together with this card's chapter etale-cohomology.\par\medskip
\noindent\textit{Editorial note.} The preamble is the chapter's own native preamble excerpt already used for this source: the Xy-pic setup (native preamble lines 14-18) and the native theorem declarations with the mathematical macros (native preamble lines 154-217).\par\medskip
\noindent\textit{Allowed standard prerequisites (editorial).}\begin{itemize}
\item Constructible sheaf and finite locally constant sheaf. Same chapter, the definitions of a constructible sheaf and of a finite locally constant sheaf on the small etale site. Supplies the three conditions compared in the statement.
\item Constructible topology on the underlying space of a quasi-compact quasi-separated scheme. Same chapter together with the chapters Topology and Properties, the fact that the constructible subsets of such a scheme form a quasi-compact boolean algebra, which turns local information into a finite partition. Supplies the passage from the local condition (2) to the partition condition (3).
\end{itemize}\medskip
\ifdefined\hypertarget\hypertarget{sourcepointer-1}{}\fi
\noindent\textit{Source reference.} Original lemma 09YD: sources/stacks-etale-cohomology/pinned-original-source.tex, source lines 11742--11742. Navigation pointer invoked at proof line 11742; the counted span contains the pointer only and the cited result is not reproduced inside it, so the pointer is registered as a non-load-bearing navigation stub in the established form for Stacks material. No bibliography citation is involved.\par\medskip
\ifdefined\hypertarget\hypertarget{sourcepointer-2}{}\fi
\noindent\textit{Source reference.} Original lemma 09Y4: sources/stacks-etale-cohomology/pinned-original-source.tex, source lines 11745--11745. Navigation pointer invoked at proof line 11745; the counted span contains the pointer only and the cited result is not reproduced inside it, so the pointer is registered as a non-load-bearing navigation stub in the established form for Stacks material. No bibliography citation is involved.\par\medskip
\ifdefined\hypertarget\hypertarget{sourcepointer-3}{}\fi
\noindent\textit{Source reference.} Original lemma 054C: sources/stacks-etale-cohomology/pinned-original-source.tex, source lines 11753--11753. Navigation pointer invoked at proof line 11753; the counted span contains the pointer only and the cited result is not reproduced inside it, so the pointer is registered as a non-load-bearing navigation stub in the established form for Stacks material. No bibliography citation is involved.\par\medskip
\ifdefined\hypertarget\hypertarget{sourcepointer-4}{}\fi
\noindent\textit{Source reference.} Original definition 03RW: sources/stacks-etale-cohomology/pinned-original-source.tex, source lines 11672--11672. Navigation pointer invoked at proof line 11755; the counted span contains the pointer only and the cited result is not reproduced inside it, so the pointer is registered as a non-load-bearing navigation stub in the established form for Stacks material. No bibliography citation is involved.\par\medskip
\begin{lemma}
\label{lemma-constructible-quasi-compact-quasi-separated}
Let $X$ be a quasi-compact and quasi-separated scheme. Let $\mathcal{F}$
be a sheaf of sets on $X_\etale$. The following are equivalent
\begin{enumerate}
\item $\mathcal{F}$ is constructible,
\item there exists an open covering $X = \bigcup U_i$ such that
$\mathcal{F}|_{U_i}$ is constructible, and
\item there exists a partition $X = \bigcup X_i$ by constructible
locally closed subschemes such that $\mathcal{F}|_{X_i}$ is finite
locally constant.
\end{enumerate}
A similar statement holds for abelian sheaves and sheaves of
$\Lambda$-modules if $\Lambda$ is Noetherian.
\end{lemma}

\begin{proof}
It is clear that (1) implies (2).

\medskip\noindent
Assume (2). For every $x \in X$ we can find an $i$ and an affine open
neighbourhood $V_x \subset U_i$ of $x$. Hence we can find a finite
affine open covering $X = \bigcup V_j$ such that for each $j$ there
exists a finite decomposition $V_j = \coprod V_{j, k}$ by locally closed
constructible subsets such that $\mathcal{F}|_{V_{j, k}}$ is finite
locally constant. By Topology, Lemma
\ref{topology-lemma-quasi-compact-open-immersion-constructible-image}
each $V_{j, k}$ is constructible as a subset of $X$.
By Topology, Lemma
\ref{topology-lemma-constructible-partition-refined-by-stratification}
we can find a finite stratification $X = \coprod X_l$ with constructible
locally closed strata such that each
$V_{j, k}$ is a union of $X_l$. Thus (3) holds.

\medskip\noindent
Assume (3) holds. Let $U \subset X$ be an affine open.
Then $U \cap X_i$ is a constructible locally closed subset of $U$
(for example by Properties, Lemma \ref{properties-lemma-locally-constructible})
and $U = \coprod U \cap X_i$ is a partition of $U$ as in
Definition \ref{definition-constructible}. Thus (1) holds.
\end{proof}

\end{document}
